% NOTES ONLY, to be written. Trimmed 2026-11-07 from thesis/methodology/04_model_size.tex
% (full draft in git, commit e4eb857).
% Outcome-blind: the budget uses only participants and events, never geometry; the selection uses
% only the literature and measurement properties, never association with plaque.
% Numbers: CGPS_data/who_where_readiness.py still uses 10/20 events-per-parameter rules; switch it to
% the Riley criteria (pmsampsize) first. Plan: ~/Desktop/Speciale/CGPS_statistics-/CLAUDE.md, open
% decision 11 (about 15 to 20 parameters at C = 0.70).
%
% The section should contain:
% - Opening: model size and features fixed before baseline geometry is linked to follow-up plaque,
%   so the LR test tests geometry, not earlier choices. First the budget, then the selection.
% - Subsection Parameter Budget (sec:method-size-budget): Riley criteria (sec:tech-sample-size),
%   shrinkage S = 0.9, anticipated C = 0.70 (source needed, e.g. PARADIGM clinical-only 0.68), also
%   0.65 and 0.75. Participants, events and P_max to fill in. Open: the criteria are for a binary
%   person-level outcome; state that n and events are counted per participant.
%   Degrees-of-freedom counting (linear 1, spline k-1, factor m-1, interaction >= 1), base first,
%   geometry gets the rest, all geometric terms linear, part of the budget left unused on purpose.
%   Table tab:method-size-budget (participants, events, P_max, base df, geometry 5, used, unused).
%   Open: SCORE2 factor coding; whether an absent-left-main indicator joins the geometric block.
% - Subsection Selection of Geometric Features (sec:method-selection): candidates are geometry the
%   literature (ch:literature) ties to WSS or plaque location and the tracer output allows. One
%   primary definition each, others to sensitivity (sec:method-validity-definitions). Candidates set
%   aside for weaker literature support (tortuosity index, bend counts, side-branch bifurcations),
%   [2026-11-07 CHECK: the Literature Review now supports tortuosity index (sec:lit-wss Kashyap,
%   sec:lit-measure), bend counts (sec:lit-measure Li) and LM-daughter angles (sec:lit-wss Cui, moved 2026-11-07);
%   "side-branch bifurcations" has no source in the review: find one or drop it.]
%   measurement error on CCTA (torsion, eccentricity, local expansion, pointwise profiles, lumen
%   volume to myocardial mass) or overlap with a retained feature (curvature x
%   calibre, upstream curvature, distance from ostium, side-branch count). Prose only, no table.
%   Four features remain: segment curvature and taper, LM angle, daughter asymmetry (2026-11-13:
%   Finet's ratio dropped, no supporting literature and the only feature needing the LM diameter;
%   see 05c header). Coverage [..]% of the cohort; the asymmetry is available whenever the angle is.
%   Budget: geometry block 4 df, plus 1 for the secondary angle x
%   asymmetry term outside the LR test. No ranking, screening, stepwise or penalised
%   selection; the rest go to secondary analyses.
% HONESTY CHECK: say the four were prioritised on these criteria, NOT that a longer list was ranked.
% 2026-11-09: the shared reason for the five (shape disturbs flow and lowers WSS, one sentence per kind
%   of feature, tied to low WSS in sec:lit-wss) is now stated once in 05_features.tex
%   (sec:method-features). sec:method-selection should only say why these four were kept over the
%   other candidates, and point to sec:method-features for the haemodynamic rationale.
\section{Parameter Budget and Feature Selection}
\label{sec:method-size}

\todonote{To be written: how many parameters the cohort supports and which geometric features fill them. See the notes in the source.}

\subsection{Parameter Budget}
\label{sec:method-size-budget}

\todonote{To be written: Riley criteria, the anticipated C-statistic and the allocation table.}

\subsection{Selection of Geometric Features}
\label{sec:method-selection}

\todonote{To be written: the candidates, why each set-aside one was set aside, and the four retained features.}
